\section{Capstone synthesis: two spectra, parity, and descriptive closure}
\label{sec:capstone-synthesis}

\subsection{Why the final synthesis is now forced}

By the end of Section~\ref{sec:meta-closure} no same-scope escape route remains open. The kernel floor is fixed, state-transition form and irreversible commitment are forced, the AMetric boundary is non-parameterized, the admissible interior is unique, mathematics is closed by operator exhaustion, physics is closed by the same standing-preserving invariant, and same-domain incompleteness objections are non-instantiable inside that interior. What remains is to classify the relation between mathematics and physics after closure and then state what that classification implies for the totality of admissible description.

The synthesis proved here is downstream and classificatory: it draws only on the already-fixed closure results for mathematics, physics, and same-domain meta-escape, together with the synthesis theorems of \eademdoc{}, and does not reopen any earlier foundational burden.\cite{MaleyEadem,MaleyClausura,MaleyFoundationalClosure,MaleyAMetric}

\subsection{Minimal background for the synthesis}

\begin{definition}[Grounding]
Within the capstone, \emph{grounding} is standing-generative authority: a domain grounds another only if it can confer standing-bearing force, foundational privilege, or admissibility-bearing status on that other domain rather than merely organize, translate, or redescribe it.
\end{definition}

\begin{definition}[Fixation mode]
A post-closure domain is \emph{internally fixed} when its admissible objects and transformations are fixed symbolically within the admitted interior itself; it is \emph{externally fixed} when its admissible use remains tied to externally fixed reports, records, or evidential anchors. These are fixation modes, not different standing values.
\end{definition}

\begin{definition}[Representational spectrum]
The \emph{representational spectrum} of a post-closure domain is the way one and the same non-generative admissible order is articulated under a fixation mode. Mathematics occupies the internally fixed symbolic spectrum; physics occupies the externally fixed empirical spectrum.
\end{definition}

\begin{definition}[Shared invariant]
The \emph{shared invariant} of the capstone is the theorem-local fact already fixed by Sections~\ref{sec:operator-exhaustion}, \ref{sec:physics-closure}, and \ref{sec:meta-closure}: standing cannot be generated by formal operators, licensed metric description, external coding, semantic import, infinitary rule-power, or same-scope extension. What survives in either domain is standing-preserving coordination inside an already fixed admissible interior.
\end{definition}

\subsection{The fixation split and two-spectrum exhaustion}

\begin{theorem}[Fixation dichotomy / two-spectrum exhaustion]
\label{thm:fixation-dichotomy}
After the closure theorems of Sections~\ref{sec:operator-exhaustion}, \ref{sec:physics-closure}, and \ref{sec:meta-closure} are fixed, every standing-preserving post-closure descriptive discipline falls into exactly one of two spectra:
\begin{enumerate}
  \item the \emph{internally fixed symbolic spectrum}, or
  \item the \emph{externally fixed empirical spectrum}.
\end{enumerate}
No third same-level admissible spectrum remains.
\end{theorem}

\begin{proof}
A post-closure discipline must still belong to the target class fixed in Sections~2--4: it must support determinate reference, admissible redescription, and standing-preserving use. If its objects and constraints are fixed symbolically from within the admitted interior itself, then it belongs to the internally fixed spectrum. If instead its evidential anchors remain externally fixed --- as reports, records, measurements, or observation-bearing carriers --- while its descriptive work remains licensed interior bookkeeping, then it belongs to the externally fixed spectrum.

No third same-level spectrum remains. A purported third spectrum would have to differ from both existing spectra by a standing-bearing invariant other than fixation mode. But Section~\ref{sec:ametric-boundary} already eliminated same-scope boundary selectors and admissibility-relevant parameterization, Section~\ref{sec:standing-regime} fixed the unique standing quotient and excluded any second same-scope standing-bearing equivalence, and Section~\ref{sec:meta-closure} excluded same-domain code, semantic, infinitary, and extension escape routes. So any purported third spectrum is either bookkeeping-only, a notational redescription of one of the two spectra, or an actual scope change rather than a third post-closure spectrum. Therefore the split is exhaustive.
\end{proof}

\begin{corollary}[Fixation is classificatory rather than hierarchical]
\label{cor:fixation-classificatory}
The distinction between the two spectra classifies how post-closure work is fixed and articulated. It does not rank one spectrum above the other in standing-bearing force.
\end{corollary}

\begin{proof}
The proof of Theorem~\ref{thm:fixation-dichotomy} introduced only a difference in fixation mode and representational articulation. No new standing-bearing source, gate, or admissibility-bearing asymmetry was added. So the split classifies the surviving difference between the two domains; it does not create a hierarchy of standing.
\end{proof}

\subsection{Closure parity}

\begin{theorem}[Closure parity theorem]
\label{thm:closure-parity}
After closure, mathematics and physics are equal in foundational status, equal in grounding incapacity, and equal in legitimate post-closure role.
\end{theorem}

\begin{proof}
Section~\ref{sec:operator-exhaustion} proved that mathematics is closed by exhaustion: no same-scope operator family generates new standing-bearing mathematical content, and what survives is internally fixed standing-preserving structural coordination. Section~\ref{sec:physics-closure} proved the parallel result for physics: no same-scope physical description generates standing, and what survives is externally fixed standing-preserving empirical coordination. Section~\ref{sec:meta-closure} then removed the only same-domain meta-level routes by which one domain might still claim standing-bearing privilege from above.

So the two domains are already equal in foundational status: neither remains foundational in a standing-generative sense. They are equal in grounding incapacity because the same shared invariant blocks mathematics from generating standing formally and blocks physics from generating standing descriptively. Finally, they are equal in legitimate post-closure role, because in each case what survives is disciplined coordination, refinement, mapping, and engineering within an already fixed admissible envelope. What differs is fixation mode and spectrum, not the standing-bearing type of role each domain is permitted to perform.
\end{proof}

\subsection{From parity to mirror classification}

\begin{proposition}[No residual distinguishing invariant]
\label{prop:no-residual-invariant}
Once foundational status, grounding incapacity, and legitimate post-closure role are fixed as in Theorem~\ref{thm:closure-parity}, no admissible invariant distinguishes mathematics and physics beyond fixation mode and the corresponding representational spectrum.
\end{proposition}

\begin{proof}
By Theorem~\ref{thm:closure-parity}, mathematics and physics already coincide in every standing-relevant classificatory dimension carried by the capstone: foundational status, grounding incapacity, and legitimate post-closure function. By Theorem~\ref{thm:fixation-dichotomy}, the only same-level admissible distinction left available is fixation mode and the resulting spectrum. Therefore any purported further invariant must be one of three things: bookkeeping that leaves the standing profile unchanged, a notational redescription of one of the already admitted spectra, or a scope change that alters the standing-bearing conditions under which the domain is classified. None of these is a residual same-level invariant beyond fixation mode.
\end{proof}

\begin{theorem}[Mirror classification theorem]
\label{thm:mirror-classification}
Mathematics and physics are equivalent under the equivalence relation induced by admissible standing-preserving structure and differ only by fixation mode and representational spectrum.
\end{theorem}

\begin{proof}
Theorem~\ref{thm:closure-parity} fixes equality of foundational status, grounding incapacity, and legitimate post-closure role. Theorem~\ref{thm:fixation-dichotomy} fixes the only surviving admissible distinction between the two domains: mathematics occupies the internally fixed symbolic spectrum and physics occupies the externally fixed empirical spectrum. Proposition~\ref{prop:no-residual-invariant} removes any further admissible invariant beyond that split. So the two domains are equivalent under admissible standing-preserving structure and differ only by fixation mode and the representational spectrum through which the same non-generative admissible order is articulated.
\end{proof}

\begin{corollary}[\emph{Eadem natura, duo spectra}]
\label{cor:eadem}
After closure, mathematics and physics are two spectra of one constrained admissible order.
\end{corollary}

\begin{proof}
Immediate from Theorem~\ref{thm:mirror-classification}. The corollary states in compressed form the same mirror-classification result: sameness in post-closure nature, difference only in spectrum.
\end{proof}

\subsection{The prohibition package}

\begin{theorem}[No grounding asymmetry]
\label{thm:no-grounding-asymmetry}
After closure, mathematics does not ground physics and physics does not ground mathematics.
\end{theorem}

\begin{proof}
Grounding in the capstone means standing-generative authority. But Theorem~\ref{thm:closure-parity} states that mathematics and physics are equal in grounding incapacity. If mathematics grounded physics, then mathematics would possess a standing-generative role denied to physics; if physics grounded mathematics, the reverse would hold. Either claim contradicts parity. Corollary~\ref{cor:fixation-classificatory} confirms that fixation mode introduces no standing-bearing hierarchy that could rescue either asymmetry.
\end{proof}

\begin{theorem}[No cross-domain rescue]
\label{thm:no-cross-domain-rescue}
The non-grounding status of mathematics cannot be repaired by appeal to physics and the non-grounding status of physics cannot be repaired by appeal to mathematics.
\end{theorem}

\begin{proof}
Suppose first that physics could rescue mathematics. Then mathematics would acquire standing-bearing force from relation to an externally fixed domain. But Section~\ref{sec:standing-regime} already excluded same-scope transfer and carrier laundering, and Theorem~\ref{thm:mirror-classification} states that the two domains differ only by fixation mode rather than by standing-bearing kind. So physics cannot rescue mathematics without smuggling in transferred standing. The converse case is symmetric: if mathematics could rescue physics by formal exactness, it would reinstall the grounding asymmetry already ruled out by Theorem~\ref{thm:no-grounding-asymmetry}. Therefore neither domain can repair the other's non-grounding status.
\end{proof}

\begin{theorem}[No meta-foundational re-escalation]
\label{thm:no-meta-reescalation}
No higher mixed standpoint, meta-foundational layer, or combined mathematics--physics platform restores foundational privilege to either domain after closure.
\end{theorem}

\begin{proof}
A purported meta-foundational re-escalation must either leave the standing profile of both domains unchanged or alter it. If it leaves the standing profile unchanged, then it is a bookkeeping redescription of the already closed relation between the two domains. If it alters that profile, then at least one domain must acquire standing-bearing authority not already available under Theorems~\ref{thm:mathematics-closure}, \ref{thm:physics-closure}, and \ref{thm:godel-non-instantiability}. But every such route would have to proceed by generation, transfer, repair, selector import, or scope change, and those routes are already exhausted by earlier sections. So no meta-foundational re-escalation remains.
\end{proof}

\begin{corollary}[No surviving reduction thesis]
\label{cor:no-reduction-thesis}
After closure, neither reduction of physics to mathematics nor reduction of mathematics to physics survives as a standing-relevant thesis.
\end{corollary}

\begin{proof}
A standing-relevant reduction thesis would have to do more than provide translation or representational economy; it would have to make one domain derivative in standing-bearing force with respect to the other. Theorems~\ref{thm:no-grounding-asymmetry} and \ref{thm:no-cross-domain-rescue} exclude exactly that. What may remain are translations, embeddings, or reformulations, but these are bookkeeping relations within the mirror structure, not reductions that reinstall privilege.
\end{proof}

\begin{corollary}[Prohibition package]
\label{cor:prohibition-package}
The following post-closure moves are inadmissible:
\begin{enumerate}
  \item re-grounding physics via mathematics;
  \item re-grounding mathematics via physics;
  \item cross-domain rescue attempts;
  \item meta-foundational re-escalation; and
  \item any standing-relevant reduction thesis between the two domains.
\end{enumerate}
\end{corollary}

\begin{proof}
Immediate from Theorems~\ref{thm:no-grounding-asymmetry}, \ref{thm:no-cross-domain-rescue}, and \ref{thm:no-meta-reescalation} together with Corollary~\ref{cor:no-reduction-thesis}.
\end{proof}

\subsection{Post-closure practice}

\begin{theorem}[Post-closure practice rule]
\label{thm:post-closure-practice}
After closure, legitimate work in mathematics and physics is restricted to standing-preserving internal transformation, representational refinement, constraint mapping, and engineering within already established admissible envelopes.
\end{theorem}

\begin{proof}
Theorem~\ref{thm:mirror-classification} identifies mathematics and physics as two spectra of one constrained admissible order. Corollary~\ref{cor:prohibition-package} removes every route by which either domain could recover standing-generative privilege over the other or over the admissible descriptive totality. So any legitimate post-closure work must remain within the standing profile already fixed for each domain. For mathematics, Section~\ref{sec:operator-exhaustion} identified the surviving role as internally fixed structural coordination. For physics, Section~\ref{sec:physics-closure} identified the surviving role as externally fixed empirical coordination. In each case the legitimate residue is therefore transformation, refinement, mapping, and engineering within the admitted interior, not new grounding.
\end{proof}

\begin{corollary}[Closure is not nihilism]
\label{cor:not-nihilism}
Closure eliminates standing-generative ambition, not substantive practice. Mathematics and physics remain fully legitimate as disciplined domains of work after closure.
\end{corollary}

\begin{proof}
Immediate from Theorem~\ref{thm:post-closure-practice}. The theorem restricts the type of legitimate work that remains; it does not eliminate that work.
\end{proof}

\begin{remark}[Seriousness without privilege]
The capstone therefore ends positively rather than nihilistically. Mathematics remains serious because exact symbolic articulation, proof, mapping, and engineering remain legitimate inside the internally fixed spectrum. Physics remains serious because measurement, modeling, empirically anchored refinement, and engineering remain legitimate inside the externally fixed spectrum. What closure removes is only the claim that either domain stands beneath the other, or beneath reality, as ground.
\end{remark}

\subsection{Descriptive closure of admissible description}

\begin{definition}[Admissible descriptive totality]
\label{def:admissible-descriptive-totality}
For the purposes of the capstone, the relevant totality is the totality of admissible descriptions under standing preservation. This is a classificatory notion, not a free-standing ontological posit over and above the fixed admissible order. The capstone therefore states closure only for admissible description, not for ontology as such.
\end{definition}

\begin{theorem}[Descriptive-closure theorem]
\label{thm:reality-closure}
Reality, in the sense of Definition~\ref{def:admissible-descriptive-totality}, is closed by exhaustion. Equivalently, once the theorem chain of this capstone is in place, no admissible description lies outside the two spectra fixed by Theorem~\ref{thm:fixation-dichotomy}, and neither spectrum contains any standing-generative route by which foundational privilege could be reintroduced.
\end{theorem}

\begin{proof}
Sections~\ref{sec:operator-exhaustion}, \ref{sec:physics-closure}, and \ref{sec:meta-closure} already closed the three substantive re-entry routes: mathematics cannot reopen the interior by operator means, physics cannot reopen it by descriptive or measurement-side means, and same-domain meta-level coding, semantic import, infinitary escalation, or extension cannot reopen it from above. Theorem~\ref{thm:fixation-dichotomy} then exhausts the remaining admissible descriptive universe into exactly two post-closure spectra. Theorem~\ref{thm:mirror-classification} shows that these are not two rival foundations but two spectra of one constrained admissible order, and Corollary~\ref{cor:prohibition-package} rules out every remaining route by which either spectrum could recover standing-generative privilege.

Therefore the totality of admissible description is closed: there is no third spectrum, no residual grounding asymmetry, no cross-domain rescue, and no same-scope meta-level reopening route. This is exactly descriptive closure in the classificatory sense fixed above. The theorem therefore closes the space of admissible description, not ontology independent of admissible description.
\end{proof}

